\section{Rigidity of the final limit}\label{sec:rigidity}

The propagation result produces contacts at shrinking scales near one terminal point. The finite slope at that point makes the final limit have constant forward and backward envelope slopes at every length. We first show why this forces an affine function.

\begin{lemma}[Rigidity from constant envelope slopes]\label{lem:slope-rigidity}
Let $v\in\Aclass_{m,C_0}$, with the class and envelopes defined in Section~\ref{sec:limiting-tools}. Suppose that for some $(b,t_0)\in\R^m\times\R$ and $B\in\R$,
\begin{equation}\label{eq:constant-envelope-slopes}
  \Senv_\ell^+v(b,t_0)=\Senv_\ell^-v(b,t_0)=B
       \qquad\text{for every }\ell>0.
\end{equation}
Then $v$ is affine.
\end{lemma}

\begin{proof}
Translate $(b,t_0)$ to $(0,0)$ and subtract the value there; write $V$ for the resulting function. These changes preserve the class and the envelope slopes. The forward equality in \eqref{eq:constant-envelope-slopes} gives $\Penv_tV(0,0)=Bt$ for $t>0$, and therefore
\begin{equation}\label{eq:future-paraboloid}
  V(y,t)\le Bt+\frac{|y|^2}{2t}\qquad(t>0).
\end{equation}

Fix any $(x,t)\in\R^m\times\R$. At a large length $\ell$, a maximizing position $y_\ell$ for $\Penv_\ell V(x,t)$ satisfies
\begin{equation}\label{eq:rigidity-optimizer-bound}
  |y_\ell|\le C_{x,t}\ell(1+\log\ell)^2.
\end{equation}
This is the optimizer estimate from Section~\ref{sec:limiting-tools}, applied on a box of size comparable to $\ell$; the translated function has the same growth estimates with constants depending on a fixed value and fit. Explicitly, if the displacement is $b_\ell\ell$ with $b_\ell\ge2$, comparison with zero displacement and the bound $|V(y,s)|\le C K(1+\log K)^2$ on boxes of size $K$ give
\[
  b_\ell\le C_{x,t}(1+\log\ell+\log b_\ell)^2.
\]
Absorbing the logarithmic term in $b_\ell$ gives $b_\ell\le C_{x,t}'(1+\log\ell)^2$, which proves \eqref{eq:rigidity-optimizer-bound}.

For $\ell$ sufficiently large, $t+\ell>0$, so \eqref{eq:future-paraboloid} applies at the optimizer. The difference between the two quadratic costs is
\begin{align*}
  \frac{|y_\ell|^2}{2(t+\ell)}-\frac{|y_\ell-x|^2}{2\ell}
  &=-\frac{t|y_\ell|^2}{2\ell(t+\ell)}
       +\frac{x\cdot y_\ell}{\ell}-\frac{|x|^2}{2\ell}\\
  &=O_{x,t}\bigl((1+\log\ell)^4\bigr)=o(\ell)
\end{align*}
by \eqref{eq:rigidity-optimizer-bound}. Consequently
\[
  \limsup_{\ell\to\infty}\Senv_\ell^+V(x,t)\le B.
\]
The forward slopes are nondecreasing in $\ell$ by Lemma~\ref{lem:envelope-slopes}. Thus $\Senv_\ell^+V(x,t)\le B$ for every $\ell>0$. Evaluating the envelope at any competitor gives the global directed comparison
\begin{equation}\label{eq:directed-comparison}
  V(y,t+\ell)-V(x,t)\le B\ell+\frac{|y-x|^2}{2\ell}
       \qquad(x,y\in\R^m,\ t\in\R,\ \ell>0).
\end{equation}

The envelopes are attained by Lemma~\ref{lem:quadratic-envelopes}. At the origin, \eqref{eq:constant-envelope-slopes} therefore supplies, for every $\ell,r>0$, a forward equality point $(\ell q_+,\ell)$ and a backward equality point $(-r q_-,-r)$, with
\[
  V(\ell q_+,\ell)=B\ell+\tfrac12\ell|q_+|^2,
  \qquad V(-r q_-,-r)=-Br-\tfrac12r|q_-|^2.
\]
Apply \eqref{eq:directed-comparison} directly from the backward point to the forward point. After canceling $B(\ell+r)$, it gives
\[
  0\ge \frac{\ell|q_+|^2+r|q_-|^2}{2}
       -\frac{|\ell q_++r q_-|^2}{2(\ell+r)}
    =\frac{\ell r}{2(\ell+r)}|q_+-q_-|^2.
\]
Hence every forward and backward equality velocity is one vector $q$, independent of all lengths and choices. We have
\[
  V(qs,s)=(B+|q|^2/2)s\qquad(s\in\R).
\]
Let $A(x,t)=q\cdot x+(B-|q|^2/2)t$. Comparing $(x,t)$ to a point $(qs,s)$ on this line in \eqref{eq:directed-comparison} yields
\[
  V(x,t)\le A(x,t)+\frac{|x-qt|^2}{2(t-s)}\quad(s<t),
  \qquad
  V(x,t)\ge A(x,t)-\frac{|x-qt|^2}{2(s-t)}\quad(s>t).
\]
Letting $s\to-\infty$ and $s\to+\infty$ proves $V=A$.
\end{proof}

\begin{proof}[Proof of Theorem~\ref{thm:liouville}]
Suppose that the class contains a nonaffine function. Proposition~\ref{prop:envelope-contact} supplies a centered solution $v$, a number $\kappa>0$, and its subsolution envelope
\[
  W(x,t)=\Cenv_t v(x,t)-\kappa^2t/2\qquad(t>0)
\]
with the stated contact and finite-slope properties. Proposition~\ref{prop:contact-propagation} supplies a point $x_{\mathrm{end}}$, infinitely many indices $j$ with $r_j=2^{-j}$, contacts $(x_j,r_j)$, and fits $(p_j,a_j)$ there of radius $r_j$, such that
\begin{equation}\label{eq:final-contact-data}
  H_v(x_{\mathrm{end}},0)>-\infty,\qquad
  W(x_j,r_j)=v(x_j,r_j),\qquad
  |x_{\mathrm{end}}-x_j+p_jr_j|\le K r_j
\end{equation}
for one finite $K$.

Apply Lemma~\ref{lem:fitted-frame-covariance} at each of these marks, and denote the transformed functions by $\bar v_j$. Their images of $(x_{\mathrm{end}},0)$ are $(b_j,0)$, where
\[
  b_j=\frac{x_{\mathrm{end}}-x_j+p_jr_j}{r_j}.
\]
These vectors are bounded by \eqref{eq:final-contact-data}. Every $\bar v_j$ has value zero and the zero fit of radius one at $(0,1)$. Lemma~\ref{lem:fit-compactness} gives a subsequence for which $\bar v_j\to v_\infty$ locally uniformly and $b_j\to b$. The uniform growth bounds also give convergence of the forward and backward envelopes at each fixed positive length, by Lemma~\ref{lem:quadratic-envelopes}.

The slope covariance \eqref{eq:fitted-slope-covariance} gives
\[
  \Senv_\ell^\pm\bar v_j(b_j,0)
    =\Senv_{r_j\ell}^\pm v(x_{\mathrm{end}},0)-a_j
       \qquad(\ell>0).
\]
Let $B=\Senv_1^+v_\infty(b,0)$, which is finite. For each fixed $\ell>0$, subtract the equality with the forward slope at length one:
\begin{align*}
  \Senv_\ell^\pm\bar v_j(b_j,0)-\Senv_1^+\bar v_j(b_j,0)
   &=\Senv_{r_j\ell}^\pm v(x_{\mathrm{end}},0)
        -\Senv_{r_j}^+v(x_{\mathrm{end}},0)\\
   &\longrightarrow0.
\end{align*}
The limit is zero because both unshifted slopes tend to the same finite value $H_v(x_{\mathrm{end}},0)$. Passing to the limit on the left gives
\[
  \Senv_\ell^+v_\infty(b,0)=\Senv_\ell^-v_\infty(b,0)=B
       \qquad(\ell>0).
\]
This argument uses each fixed length along the same locally convergent subsequence; no convergence of the individual shifts $a_j$ is required. Lemma~\ref{lem:slope-rigidity} now makes $v_\infty$ affine.

The contacts retain a positive envelope gap. By \eqref{eq:contact-barycentric-gap} at the contacts in \eqref{eq:final-contact-data} and the gap covariance \eqref{eq:contact-W-covariance},
\[
  \Cenv_1\bar v_j(0,1)-\bar v_j(0,1)=\kappa^2/2.
\]
Lemma~\ref{lem:barycentric-envelope} passes this identity to the locally uniform limit with its uniform growth bound. Thus
\[
  \Cenv_1v_\infty(0,1)-v_\infty(0,1)=\kappa^2/2>0.
\]
For an affine function, the average of its values at any configuration is its value at the barycenter, so the nonnegative variance penalty makes $\Cenv_1v_\infty=v_\infty$. This contradiction proves Theorem~\ref{thm:liouville}. The reduction in Section~\ref{sec:reduction} then proves Theorem~\ref{thm:main}.
\end{proof}
